\documentclass[10pt,a4paper]{amsart}
\usepackage{amsmath,amssymb,amsthm,graphicx}
% Retained minimal preamble for the Stacks Project chapter more-algebra.tex (More on Algebra).
% The chapter file pinned-original-source.tex carries only \input{preamble} on its first line, so no
% preamble text lives inside the chapter itself.  The Stacks Project uses one single shared preamble for
% every chapter; the snapshot stages it identically as sources/stacks-chow/preamble.tex and
% sources/stacks-etale-cohomology/preamble.tex.  The block below is copied mechanically, line for line,
% from sources/stacks-chow/preamble.tex lines 14-18 (the Xy-pic setup) and 154-217 (the native theorem
% declarations together with the mathematical macros).  Nothing was added, removed or redefined; the
% wrapper class and the amsmath/amssymb/amsthm/graphicx base packages are supplied by the assembler.
% Chapter: more-algebra.tex -- More on Algebra.
% For commutative diagrams we use Xy-pic
\usepackage[all]{xy}

% We use 2cell for 2-commutative diagrams.
\xyoption{2cell}

% Theorem environments.
%
\theoremstyle{plain}
\newtheorem{theorem}[subsection]{Theorem}
\newtheorem{proposition}[subsection]{Proposition}
\newtheorem{lemma}[subsection]{Lemma}

\theoremstyle{definition}
\newtheorem{definition}[subsection]{Definition}
\newtheorem{example}[subsection]{Example}
\newtheorem{exercise}[subsection]{Exercise}
\newtheorem{situation}[subsection]{Situation}

\theoremstyle{remark}
\newtheorem{remark}[subsection]{Remark}
\newtheorem{remarks}[subsection]{Remarks}

\numberwithin{equation}{subsection}

% Macros
%
\def\lim{\mathop{\mathrm{lim}}\nolimits}
\def\colim{\mathop{\mathrm{colim}}\nolimits}
\def\Spec{\mathop{\mathrm{Spec}}}
\def\Hom{\mathop{\mathrm{Hom}}\nolimits}
\def\Ext{\mathop{\mathrm{Ext}}\nolimits}
\def\SheafHom{\mathop{\mathcal{H}\!\mathit{om}}\nolimits}
\def\SheafExt{\mathop{\mathcal{E}\!\mathit{xt}}\nolimits}
\def\Sch{\mathit{Sch}}
\def\Mor{\mathop{\mathrm{Mor}}\nolimits}
\def\Ob{\mathop{\mathrm{Ob}}\nolimits}
\def\Sh{\mathop{\mathit{Sh}}\nolimits}
\def\NL{\mathop{N\!L}\nolimits}
\def\CH{\mathop{\mathrm{CH}}\nolimits}
\def\proetale{{pro\text{-}\acute{e}tale}}
\def\etale{{\acute{e}tale}}
\def\QCoh{\mathit{QCoh}}
\def\Ker{\mathop{\mathrm{Ker}}}
\def\Im{\mathop{\mathrm{Im}}}
\def\Coker{\mathop{\mathrm{Coker}}}
\def\Coim{\mathop{\mathrm{Coim}}}

% Boxtimes
%
\DeclareMathSymbol{\boxtimes}{\mathbin}{AMSa}{"02}

%
% Macros for moduli stacks/spaces
%
\def\QCohstack{\mathcal{QC}\!\mathit{oh}}
\def\Cohstack{\mathcal{C}\!\mathit{oh}}
\def\Spacesstack{\mathcal{S}\!\mathit{paces}}
\def\Quotfunctor{\mathrm{Quot}}
\def\Hilbfunctor{\mathrm{Hilb}}
\def\Curvesstack{\mathcal{C}\!\mathit{urves}}
\def\Polarizedstack{\mathcal{P}\!\mathit{olarized}}
\def\Complexesstack{\mathcal{C}\!\mathit{omplexes}}
% \Pic is the operator that assigns to X its picard group, usage \Pic(X)
% \Picardstack_{X/B} denotes the Picard stack of X over B
% \Picardfunctor_{X/B} denotes the Picard functor of X over B
\def\Pic{\mathop{\mathrm{Pic}}\nolimits}
\def\Picardstack{\mathcal{P}\!\mathit{ic}}
\def\Picardfunctor{\mathrm{Pic}}
\def\Deformationcategory{\mathcal{D}\!\mathit{ef}}
\usepackage{cleveref}
\crefname{originallemma}{original lemma}{original lemmas}
\Crefname{originallemma}{Original lemma}{Original lemmas}
\makeatletter
\newlabel{more-algebra-lemma-fitting-ideal}{{07Z8}{source}{}{sourcepointer-1}{}}
\newlabel{more-algebra-lemma-fitting-ideal@cref}{{[originallemma][0][]07Z8}{[1][0][]source}{}{}{}}
\makeatother
\crefname{originallemma}{original lemma}{original lemmas}
\Crefname{originallemma}{Original lemma}{Original lemmas}
\makeatletter
\newlabel{lemma-ideals-generated-by-minors}{{07Z7}{source}{}{sourcepointer-2}{}}
\newlabel{lemma-ideals-generated-by-minors@cref}{{[originallemma][0][]07Z7}{[1][0][]source}{}{}{}}
\makeatother
\crefname{originallemma}{original lemma}{original lemmas}
\Crefname{originallemma}{Original lemma}{Original lemmas}
\makeatletter
\newlabel{algebra-lemma-Schanuel}{{00O3}{source}{}{sourcepointer-3}{}}
\newlabel{algebra-lemma-Schanuel@cref}{{[originallemma][0][]00O3}{[1][0][]source}{}{}{}}
\makeatother
\title{Local freeness detected by the Fitting ideal of a presentation}
\author{The Stacks Project Authors}
\date{stacks-project:more-algebra.tex:local-snapshot-20261012; GNU FDL 1.2 or later}
\begin{document}
\maketitle
\setcounter{subsection}{1}
\noindent\textit{Editorial note.} All delivered mathematical text is the author's own, copied from the pinned Stacks Project chapter file more-algebra.tex; no statement or proof step was rewritten, reordered or removed.\par\medskip
\noindent\textit{Editorial note.} The chapter file pinned-original-source.tex carries only \textbackslash{}input\{preamble\} on its first line, so no preamble text lives inside the chapter. The retained file sources/stacks-more-algebra/retained-preamble.tex is a mechanical, line-for-line copy of the project's single shared preamble as staged in the snapshot (Xy-pic setup plus theorem declarations and mathematical macros). Nothing was added, removed or redefined, and the excerpt serves only as preamble, never as a span coordinate.\par\medskip
\noindent\textit{Editorial note.} Both counted spans start at the \textbackslash{}begin line of their environment, so the statement and the author's proof are delivered whole; surrounding prose of the section is deliberately outside the counted spans and is carried as standard prerequisites instead.\par\medskip
\noindent\textit{Editorial note.} The citation fields (source\_id, source\_version, authors, source\_url, license.name, source\_title) are identical to those of every other card this lane has produced from this chapter, so the batch selector treats them as one source.\par\medskip
\noindent\textit{Allowed standard prerequisites (editorial).}\begin{itemize}
\item Fitting ideal of a finite module, defined through the minors of a presentation matrix. defined in the same section, source lines 1320-1354, before the counted statement. Supplies the ideal whose relation to local freeness is asserted.
\item Criterion for local freeness of a finite module over a ring through the rank function and the Fitting ideals. standard, used at proof lines 1380-1400. Converts the rank condition into the statement that a minor generates the unit ideal.
\item Behaviour of the maximal minors of a matrix under localisation. used in both directions of the proof, lines 1370-1410. Lets the argument be carried out after localising at a prime.
\end{itemize}\medskip
\ifdefined\hypertarget\hypertarget{sourcepointer-1}{}\fi
\noindent\textit{Source reference.} Original lemma 07Z8: focused-seeds/AGP/sources/stacks-snapshot-20261011/tags.txt, source lines 10138--10138. The lemma's own canonical tag, quoted so that readers can resolve the label inside the delivered text; it points at the counted statement itself and carries no proof content.\par\medskip
\ifdefined\hypertarget\hypertarget{sourcepointer-2}{}\fi
\noindent\textit{Source reference.} Original lemma 07Z7: focused-seeds/AGP/sources/stacks-snapshot-20261011/tags.txt, source lines 10137--10137. Navigation pointer to a statement the counted proof refers to but does not contain; nothing load-bearing is delegated to it.\par\medskip
\ifdefined\hypertarget\hypertarget{sourcepointer-3}{}\fi
\noindent\textit{Source reference.} Original lemma 00O3: focused-seeds/AGP/sources/stacks-snapshot-20261011/tags.txt, source lines 881--881. Navigation pointer to a statement the counted proof refers to but does not contain; nothing load-bearing is delegated to it.\par\medskip
\begin{lemma}
\label{lemma-fitting-ideal}
Let $R$ be a ring. Let $M$ be a finite $R$-module. Choose a presentation
$$
\bigoplus\nolimits_{j \in J} R \longrightarrow R^{\oplus n}
\longrightarrow M \longrightarrow 0.
$$
of $M$. Let $A = (a_{ij})_{i = 1, \ldots, n, j \in J}$ be the matrix
of the map $\bigoplus_{j \in J} R \to R^{\oplus n}$.
The ideal $\text{Fit}_k(M)$ generated by the
$(n - k) \times (n - k)$ minors of
$A$ is independent of the choice of the presentation.
\end{lemma}

\begin{proof}
Let $K \subset R^{\oplus n}$ be the kernel of the surjection
$R^{\oplus n} \to M$. Pick $z_1, \ldots, z_{n - k} \in K$
and write $z_j = (z_{1j}, \ldots, z_{nj})$.
Another description of the ideal $\text{Fit}_k(M)$
is that it is the ideal generated by the $(n - k) \times (n - k)$ minors of
all the matrices $(z_{ij})$ we obtain in this way.

\medskip\noindent
Suppose we change the surjection into the surjection
$R^{\oplus n + n'} \to M$ with kernel $K'$ where we use the original
map on the first $n$ standard basis elements of $R^{\oplus n + n'}$
and $0$ on the last $n'$ basis vectors. Then the corresponding ideals
are the same. Namely, if $z_1, \ldots, z_{n - k} \in K$ as above,
let $z'_j = (z_{1j}, \ldots, z_{nj}, 0, \ldots, 0) \in K'$ for
$j = 1, \ldots, n - k$ and
$z'_{n + j'} = (0, \ldots, 0, 1, 0, \ldots, 0) \in K'$. Then we see that
the ideal of $(n - k) \times (n - k)$ minors of $(z_{ij})$ agrees
with the ideal of $(n + n' - k) \times (n + n' - k)$ minors of
$(z'_{ij})$. This gives one of the inclusions.
Conversely, given  $z'_1, \ldots, z'_{n + n' - k}$
in $K'$ we can project these to $R^{\oplus n}$ to get
$z_1, \ldots, z_{n + n' - k}$ in $K$. By
Lemma \ref{lemma-ideals-generated-by-minors}
we see that the ideal generated by the
$(n + n' - k) \times (n + n' - k)$ minors of
$(z'_{ij})$ is contained in the ideal generated by the
$(n - k) \times (n - k)$ minors of $(z_{ij})$. This gives the
other inclusion.

\medskip\noindent
Let $R^{\oplus m} \to M$ be another surjection with kernel $L$.
By Schanuel's lemma (Algebra, Lemma \ref{algebra-lemma-Schanuel})
and the results of the previous paragraph, we may assume $m = n$
and that there is an isomorphism $R^{\oplus n} \to R^{\oplus m}$
commuting with the surjections to $M$. Let $C = (c_{li})$ be the
(invertible) matrix of this map (it is a square matrix as $n = m$).
Then given $z'_1, \ldots, z'_{n - k} \in L$ as above we can find
$z_1, \ldots, z_{n - k} \in K$ with
$z_1' = Cz_1, \ldots, z'_{n - k} = Cz_{n - k}$. By
Lemma \ref{lemma-ideals-generated-by-minors} we get one of the
inclusions. By symmetry we get the other.
\end{proof}

\end{document}
